\documentclass[UTF8,a4paper,11pt]{ctexart}

\usepackage{amsmath,amssymb,bm,mathtools}
\usepackage{booktabs}
\usepackage{array}
\usepackage{enumitem}
\usepackage{geometry}
\usepackage{listings}
\usepackage{xcolor}
\usepackage{xurl}
\usepackage{algorithm}
\usepackage{algpseudocode}
\usepackage[colorlinks=true,linkcolor=blue,citecolor=blue,urlcolor=blue]{hyperref}

\geometry{left=2.5cm,right=2.5cm,top=2.6cm,bottom=2.6cm}
\setlist{nosep}
\linespread{1.15}

\definecolor{codebg}{RGB}{248,248,248}
\definecolor{codekw}{RGB}{0,80,160}
\definecolor{codecm}{RGB}{0,128,0}
\definecolor{codestr}{RGB}{163,21,21}

\lstset{
  basicstyle=\ttfamily\footnotesize,
  backgroundcolor=\color{codebg},
  keywordstyle=\color{codekw}\bfseries,
  commentstyle=\color{codecm}\itshape,
  stringstyle=\color{codestr},
  language=C++,
  showstringspaces=false,
  breaklines=true,
  columns=fullflexible,
  frame=single,
  framesep=4pt,
  xleftmargin=6pt,
  xrightmargin=6pt,
}

\newcommand{\dd}{\mathrm{d}}
\newcommand{\jj}{\mathrm{j}}
\newcommand{\RR}{\mathbb{R}}
\newcommand{\CC}{\mathbb{C}}
\newcommand{\transpose}{\mathsf{T}}
\newcommand{\herm}{\mathsf{H}}
\newcommand{\Ne}{\bm{N}}
\newcommand{\Ke}{\bm{K}}
\newcommand{\Me}{\bm{M}}
\newcommand{\code}[1]{\texttt{\small #1}}

\title{\texorpdfstring{\code{FEMAssembler::buildAffineSystem()} 数学模型与实现解读}{FEMAssembler::buildAffineSystem() 数学模型与实现解读}}
\author{独立推导说明}
\date{\today}

\begin{document}

\maketitle

\begin{abstract}
\code{FEMAssembler::buildAffineSystem()} 是本工程把频域有限元 curl--curl 系统拆成
\emph{频率无关}\ 部分的入口，是 ALPS 快速扫频（\code{AlpsSweep}）以及 Krylov 类模型降阶（MOR）方法可以直接消费的标准化输入。本文以 Nédélec/层次向量有限元、微波器件三维切向矢量有限元、TFE 与 ALPS/MOR 文献为背景
\cite{Nedelec1980,Graglia1997,Webb1999,Cendes1988,Lee1990,Bracken1998,Sun2001,Anderson2004}，
从频域 Maxwell 方程和棱元 Galerkin 离散出发，推导该函数返回的矩阵、端口耦合向量和右端项的物理含义，并把推导各步骤对回到 \code{src/fem/FEMAssembler.cpp}、\code{src/bc/WavePortBC.cpp} 与 \code{src/sweep/AlpsSweep.cpp} 中的实现。
\end{abstract}

\tableofcontents

\section{为什么需要"仿射"系统}
\label{sec:why-affine}

频域电磁有限元在每个频点 $\omega$（角频率，$\omega=2\pi f$）处都需要装配并求解一个大型复对称稀疏线性系统
\begin{equation}
    \bm A(\omega)\,\bm x(\omega)=\bm b(\omega).
    \label{eq:per-freq}
\end{equation}
若每个频点都从零开始走完"算单元矩阵 $\to$ 累加全局矩阵 $\to$ 因式分解 $\to$ 回代"四步，扫频成本随频点数线性增长。把 $\bm A(\omega)$ 写成
\begin{equation}
    \bm A(\omega)=\sum_{q} f_q(\omega)\,\bm A_q
    \label{eq:affine-form}
\end{equation}
的"仿射展开"（affine decomposition）后，所有 $\bm A_q$ 都是\emph{频率无关}\ 的稀疏矩阵，
只在 $\omega$ 上做轻量加权。这正是 Krylov MOR、AWE/ALPS 等快速扫频方法
赖以工作的代数前提。\code{FEMAssembler::buildAffineSystem()} 就是一次性产出
\eqref{eq:affine-form} 中所有 $\bm A_q$ 及其加权函数所需要的物理量。

下文先在 §\ref{sec:weak-form} 推导 $\bm A(\omega)$ 的连续与离散形式，§\ref{sec:affine}
给出仿射拆分公式，§\ref{sec:port-affine} 推导端口算子的秩 1 拆分。
§\ref{sec:impl} 把这些数学量逐项映射到 C++ 实现，§\ref{sec:invariants}--\ref{sec:limits}
讨论实现层的不变量、PEC 投影、损耗判定与目前的限制。

\section{从 Maxwell 方程到 curl--curl 弱形式}
\label{sec:weak-form}

约定时间因子 $e^{\jj\omega t}$。各向同性线性介质中，频域 Maxwell 方程
\begin{align}
    \nabla\times\bm E&=-\jj\omega\mu\bm H,\\
    \nabla\times\bm H&=(\sigma+\jj\omega\varepsilon)\bm E+\bm J_s
\end{align}
消去 $\bm H$ 得电场 curl--curl 方程
\begin{equation}
    \nabla\times\bigl(\mu_r^{-1}\nabla\times\bm E\bigr)
    -k_0^{2}\,\varepsilon_c(\omega)\,\bm E
    =-\jj\omega\mu_0\bm J_s,
    \label{eq:curlcurl}
\end{equation}
其中
\begin{equation}
    k_0=\omega\sqrt{\mu_0\varepsilon_0},\qquad
    \varepsilon_c(\omega)=\varepsilon_r-\jj\,\frac{\sigma}{\omega\varepsilon_0}.
    \label{eq:eps_c}
\end{equation}

记计算域为 $\Omega$，端口面记作 $\Gamma_p\subset\partial\Omega$（$p=1,\dots,N_p$），
PEC 面记作 $\Gamma_e$。乘以测试函数 $\bm v\in H_0(\mathrm{curl};\Omega)$（在 $\Gamma_e$ 上
满足 $\bm n\times\bm v=0$）并对 $\Omega$ 积分，对 curl 项做分部积分，得到弱形式
\begin{equation}
    \int_\Omega\!\mu_r^{-1}(\nabla\times\bm v)\!\cdot\!(\nabla\times\bm E)\,\dd V
    -k_0^{2}\!\int_\Omega\!\varepsilon_c\,\bm v\!\cdot\!\bm E\,\dd V
    +\sum_{p}\!\int_{\Gamma_p}\!\bm v\!\cdot\!\bigl(\bm n\times\mu_r^{-1}\nabla\times\bm E\bigr)\,\dd S
    =\text{(右端)}.
    \label{eq:weak}
\end{equation}
最关键的一步是把 $\Gamma_p$ 上的边界项代换为只用切向电场表达的端口算子，
推导见 §\ref{sec:port-affine}。
这里的边界项符号采用外法向 $\bm n$ 与时间因子 $e^{\jj\omega t}$；若端口模的传播方向或相量约定改变，
端口算子的整体符号会随之改变。本工程的实现约定与 \code{WavePortBC::applyOneMode} 中的
\code{admittance = j beta} 一致。

\subsection{Galerkin 棱元离散}

把 $\bm E$ 在棱元（Nédélec H(curl) 一阶不完备 / 二阶半完备）基 $\{\Ne_i\}_{i=1}^{N}$
上展开
\begin{equation}
    \bm E\approx\sum_i x_i\,\Ne_i,
\end{equation}
取 $\bm v=\Ne_a$，$\bm E$ 的展开系数为 $x_b$，则弱形式 \eqref{eq:weak} 中的体积分两项
直接给出全局矩阵
\begin{equation}
    A_{ab}^{\text{vol}}(\omega)
    =\underbrace{\int_\Omega\!\mu_r^{-1}(\nabla\!\times\!\Ne_a)\!\cdot\!(\nabla\!\times\!\Ne_b)\,\dd V}_{\equiv K_{ab}}
    -k_0^2\underbrace{\int_\Omega\!\varepsilon_c\,\Ne_a\!\cdot\!\Ne_b\,\dd V}_{\equiv \tilde M_{ab}(\omega)}.
    \label{eq:Avol}
\end{equation}
注意 $K_{ab}$ 完全频率无关，而 $\tilde M_{ab}(\omega)$ 因含 $\varepsilon_c(\omega)$ 在
$\sigma\neq0$ 时仍是频率函数。

按单元 $T$ 分解，标准元素装配
\begin{equation}
    K_{ab}=\sum_T\frac{1}{\mu_{r,T}}\,\Ke^T_{ab},\qquad
    \Ke^T_{ab}=\!\int_T(\nabla\!\times\!\Ne_a)\!\cdot\!(\nabla\!\times\!\Ne_b)\,\dd V,
    \label{eq:Ke}
\end{equation}
\begin{equation}
    \tilde M_{ab}(\omega)=\sum_T\varepsilon_{c,T}(\omega)\,\Me^T_{ab},\qquad
    \Me^T_{ab}=\!\int_T\Ne_a\!\cdot\!\Ne_b\,\dd V.
    \label{eq:Me}
\end{equation}
$\Ke^T,\Me^T$ 只依赖单元几何与基函数，不依赖 $\omega$，是\emph{频率无关}\ 的
单元参考矩阵。

\section{\texorpdfstring{$\bm A(\omega)$}{A(omega)} 的仿射拆分}
\label{sec:affine}

\subsection{\texorpdfstring{无损情形：单一 $\bm M$}{无损情形：单一 M}}

无损情形 $\sigma\equiv0$，因此 $\varepsilon_c=\varepsilon_r$ 也与 $\omega$ 无关。
体积部分立刻退化为
\begin{equation}
    \bm A^{\text{vol}}(\omega)=\bm K-k_0^{2}(\omega)\,\bm M,
    \label{eq:Avol-lossless}
\end{equation}
其中
\begin{equation}
    \bm K=\sum_T\frac{1}{\mu_{r,T}}\,\bm K^T,\qquad
    \bm M=\sum_T\varepsilon_{r,T}\,\bm M^T.
    \label{eq:K-M-global}
\end{equation}
两者都是\emph{实}\ 对称稀疏矩阵，与 $\omega$ 完全解耦。把端口算子（§\ref{sec:port-affine}）
合并进来后，传播模对应的线性系统可写为
\begin{equation}
    \boxed{\;\bm A(\omega)=\bm K-k_0^{2}\,\bm M
    +\sum_{v=1}^{N_v}Y_v(\omega)\,\bm m_v\bm m_v^{\transpose}\;},
    \qquad
    \boxed{\;\bm b(\omega)=2\sum_{v=1}^{N_v}Y_v(\omega)a_v^{\rm inc}\,\bm m_v\;}
    \label{eq:affine-lossless}
\end{equation}
其中 $N_v$ 是\emph{虚拟端口}\ 数（每个物理端口的每个保留模式各一个），
$\bm m_v$ 是端口耦合稀疏向量，
\begin{equation}
    Y_v(\omega)=\jj\,\beta_v^+(\omega),\qquad
    \beta_v^+(\omega)=\sqrt{\max(0,k_0^2-k_{c,v}^2)}
    \label{eq:modal-admittance}
\end{equation}
是当前 affine/ALPS 路径实际使用的简化模导纳。$a_v^{\rm inc}$ 为入射模幅度；未被标为
\code{isExcitationMode} 的虚拟端口有 $a_v^{\rm inc}=0$。单频点 direct path 为了数值稳健，
在 $\beta_v^+=0$ 时使用 \code{j*k0} 的占位导纳；该低于截止模处理不是严格的波导模导纳，
也不是当前 ALPS reduced evaluation 的一部分，见 §\ref{sec:limits}。

\subsection{有损情形：单一质量矩阵不再充分}

若存在 $\sigma>0$，则 \eqref{eq:eps_c} 给出 $\varepsilon_c(\omega)=\varepsilon_r-\jj\sigma/(\omega\varepsilon_0)$，
\eqref{eq:Avol} 中的质量项变为
\begin{equation}
    \tilde M_{ab}(\omega)
    =\sum_T\Bigl(\varepsilon_{r,T}-\jj\,\frac{\sigma_T}{\omega\varepsilon_0}\Bigr)\Me^T_{ab}
    =\varepsilon_r\text{-加权全局} M
    -\frac{\jj}{\omega\varepsilon_0}\,(\sigma\text{-加权全局} M).
\end{equation}
这等价于把单一 $\bm M$ 拆成 \emph{两个}\ 全局参考矩阵
\begin{equation}
    \bm M=\sum_T\varepsilon_{r,T}\bm M^T,\qquad
    \bm M_\sigma=\sum_T\sigma_T\bm M^T,
\end{equation}
若显式缓存 $\bm M_\sigma$，体积项仍可写成多项仿射形式
\begin{equation}
    \bm A(\omega)=\bm K-k_0^{2}\bm M
    +\frac{\jj k_0^{2}}{\omega\varepsilon_0}\bm M_\sigma
    +\sum_vY_v(\omega)\bm m_v\bm m_v^{\transpose}.
\end{equation}
当前实现\emph{尚未}\ 缓存 $\bm M_\sigma$，所以遇到 $\sigma>0$ 时 \code{AffineSystem::lossless}
被置为 \code{false}，调用方（\code{AlpsSweep}）改走逐频点 direct path。
具体判定见 §\ref{sec:lossless-flag}。

\section{端口边界项的秩 1 仿射拆分}
\label{sec:port-affine}

\subsection{端口模的来源：默认走二维端面本征}
\label{sec:port-mode-source}

仿射拆分需要每个端口面 $\Gamma_p$ 上的"传播模式形状" $\bm e_{p,m}^{\perp}(\bm r)$
和对应截止波数 $k_{c,p,m}^{2}$。这种把端口区模函数与三维有限元未知量用变分形式耦合的思想，
与微波器件 FEM/TFE 文献中的端口模展开一致 \cite{Cendes1988,Lee1990}。本工程现支持三条路径产出这两个量，
全部最终落到同一 \code{PortMode} 数据契约上：
\begin{enumerate}[leftmargin=2em]
\item \textbf{NPM (\code{--port-method numerical}，默认)}：把端口面 $\Gamma_p$ 上的
三角剖分单独抽出来当作 2D 网，在其上做 H(curl) 棱元离散，
求解二维广义本征值问题
\begin{equation}
    \bm K_{\text{port}}\,\bm v=k_{c}^{2}\,\bm M_{\text{port}}\,\bm v,
    \label{eq:port-eig}
\end{equation}
取最低非伪本征对 $(\bm v_1,k_{c,1}^{2})$ 作为主模。模式形状由
$\bm e^{\perp}=\sum_i v_i\,\Ne_i^{\partial}$ 给出，其中 $\Ne_i^{\partial}$ 是端口面
切向 H(curl) 棱元基。对任何二维截面（矩形 / 脊形 / 圆形 / 偏心）都成立。
实现：\code{PortModeSolver::solve(faceId)} → \code{computeMode} →
\code{computeMultiMode(faceId, 1)}。
\item \textbf{APM (\code{--port-method analytic})}：直接套矩形波导 TE10 闭式
$\bm e_{\text{TE10}}=\sqrt{2/(ab)}\sin(\pi u/a)\,\hat v$，
$k_c^{2}=(\pi/a)^{2}$，再 L2 投影到 FE 端口 dof。仅矩形截面可用。
\item \textbf{TFE (\code{--port-method tfe})}：与 NPM 同一本征解
\eqref{eq:port-eig}，但保留最低 $N$ 个非伪本征对，每个模出一个独立的虚拟端口。
\end{enumerate}
NPM 和 TFE 共用 \code{PortModeSolver::computeMultiMode}：单元矩阵对应面 H(curl)
基的 curl-curl 与 mass，三角形 7 点 Wandzura 求积；本征求解走
\code{LAPACKE\_dsygv}（无 MKL 时回落 Jacobi 迭代）；
零本征对应离散梯度零空间，按相对阈值剔除。
当前本征式 \eqref{eq:port-eig} 未显式带入端口截面材料张量，因此等价于空气/均匀归一化端口；
若后续支持介质填充或色散端口，应把 $\beta^2$ 推广为 $k_0^2\mu_r\varepsilon_r-k_c^2$
或直接构造带材料权重的端面本征问题。

\subsection{从端面本征到端口算子的代数推导}
\label{sec:port-affine-derivation}

设 $\Gamma_p$ 上传播波导支持 $N_{m,p}$ 个传播或近截止模 $\bm e_{p,m}^{\perp}(\bm r)$，
$m=1,\dots,N_{m,p}$。无论这些模是来自数值本征 \eqref{eq:port-eig} 还是闭式 TE10，
当前 affine 路径只把传播部分写成
\begin{equation}
    \beta_{p,m}^+(\omega)=\sqrt{\max\bigl(0,k_0^{2}-k_{c,p,m}^{2}\bigr)},\qquad
    Y_{p,m}(\omega)=\jj\beta_{p,m}^+(\omega).
    \label{eq:beta}
\end{equation}
在 $\Gamma_p$ 上令真正的切向投影为
\begin{equation}
    \bm E_t=(\bm I-\bm n\bm n^{\transpose})\bm E
           =-\bm n\times(\bm n\times\bm E).
    \label{eq:tangential-projection}
\end{equation}
再把切向电场展开为
\begin{equation}
    \bm E_t\big|_{\Gamma_p}
    =\sum_{m=1}^{N_{m,p}} \alpha_{p,m}\,\bm e_{p,m}^{\perp}+\bm E^{\text{inc}},
    \label{eq:port-expansion}
\end{equation}
代入 \eqref{eq:weak} 中的端口边界积分后，无源端口在弱形式中贡献的部分是
\begin{equation}
    \int_{\Gamma_p}\bm v\cdot
    \Bigl[\sum_mY_{p,m}(\omega)\bigl(\bm e_{p,m}^{\perp}\cdot\bm E_t\bigr)\bm e_{p,m}^{\perp}\Bigr]\dd S,
    \label{eq:port-bilinear}
\end{equation}
这一步的来源是把端口边界上的切向场先展开，再用模正交性把边界算子投影回单个模方向。更具体地说，若将
\begin{equation}
    \bm E_t\big|_{\Gamma_p}=\sum_m \alpha_{p,m}\bm e_{p,m}^{\perp}+\bm E^{\text{inc}}
\end{equation}
代入端口边界条件，那么对第 $m$ 个传播模会得到一个标量关系
\begin{equation}
    \bm n\times\mu_r^{-1}\nabla\times\bm E
    =-\sum_m Y_{p,m}(\omega)\,\alpha_{p,m}\bm e_{p,m}^{\perp}
    +2\sum_m Y_{p,m}(\omega)\,a_{p,m}^{\rm inc}\bm e_{p,m}^{\perp}.
\end{equation}
这里第一项对应“由未知场自身反射回去的反应电流”，第二项对应“已知入射波注入的激励电流”。
把它代回 \eqref{eq:weak} 的边界积分后，得到双线性项
\begin{equation}
    \sum_m Y_{p,m}(\omega)\int_{\Gamma_p}(\bm v\cdot\bm e_{p,m}^{\perp})(\bm E_t\cdot\bm e_{p,m}^{\perp})\,\dd S,
\end{equation}
以及对应的右端项
\begin{equation}
    \bm f_{p}^{\rm port}(\omega)
    =2\sum_m Y_{p,m}(\omega)a_{p,m}^{\rm inc}
    \int_{\Gamma_p}(\bm v\cdot\bm e_{p,m}^{\perp})\,\dd S.
    \label{eq:port-rhs-weak}
\end{equation}
对有限元试函数取 \(\bm v=\Ne_a\) 后，定义端口耦合向量
\begin{equation}
    [\bm m_{p,m}]_a=\int_{\Gamma_p}\Ne_a\cdot\bm e_{p,m}^{\perp}\,\dd S,
\end{equation}
就得到矩阵形式的秩 1 算子
\begin{equation}
    \sum_mY_{p,m}(\omega)\,\bm m_{p,m}\bm m_{p,m}^{\transpose},
    \label{eq:m_p}
\end{equation}
以及右端
\begin{equation}
    \bm f_{p}^{\rm port}(\omega)=2\sum_mY_{p,m}(\omega)a_{p,m}^{\rm inc}\bm m_{p,m}.
\end{equation}
因此，公式右端中的系数“2”并不是经验因子，而是端口边界条件里入射波与反射波叠加后留下的标准结果：当把总场写成“未知散射场 + 已知入射场”时，未知场产生一个负号，入射场又额外贡献一次同样大小的项，二者合并后就得到 2。

从散射参数的角度看，这个右端项也可以这样理解：端口模展开通常把每个虚拟端口上的场写成“入射幅度 $a_{p,m}^{\rm inc}$ + 反射幅度 $a_{p,m}^{\rm ref}$”，而目标求解的并不是单独的边界场，而是给定入射激励后系统对散射波的响应。对于传播模，边界条件实际上把未知的反射波通量与端口导纳 $Y_{p,m}(\omega)$ 绑定起来，因此当把边界算子移到弱式右端时，已知的入射幅度就会以线性源项的形式出现，并且与同一个耦合向量 $\bm m_{p,m}$ 相乘。这就是为什么端口散射参数、模导纳和有限元右端项在代数上最终会写成同一个低秩结构：前者定义物理输入输出，后者则是其有限元离散后的矩阵表达。

$\bm m_{p,m}$ 在代码中称为 \code{couplingWeights}，是\emph{稀疏}\ 的——只有支撑落在
$\Gamma_p$ 的边自由度的分量非零。$\bm m_{p,m}$ 与 $\omega$ 无关；
$Y_{p,m}(\omega)$ 是端口项的频率函数。若该模式带入射幅度 $a_{p,m}^{\rm inc}$，
右端同步获得 $2Y_{p,m}a_{p,m}^{\rm inc}\bm m_{p,m}$，这正对应
\code{WavePortBC::applyOneMode} 中的 \code{rhs += 2.0 * admittance * incident * rowValue}。
NPM/TFE 路径下 $\bm m_{p,m}$ 直接由 $\bm M_{\text{port}}\bm v_m$ 给出（
$M$-内积投影），APM 路径下 $\bm m_{p,m}$ 是闭式 $\bm e^{\perp}$ 在 H(curl) 端口 dof 上
的 L2 投影。三者共用同一 \code{PortMode::couplingWeights} 字段。
\code{computeMultiMode} 会把本征向量按 $\bm v_m^{\transpose}\bm M_{\text{port}}\bm v_m=1$
归一化，再保存 $\bm m_m=\bm M_{\text{port}}\bm v_m$；因此
$\bm m_m^{\transpose}\bm x$ 就是离散场在该端口模上的 $M$-内积投影，而不是原始边 dof 系数。

\subsection{虚拟端口列表}

把 \eqref{eq:m_p} 对所有物理端口、所有解析模式平铺成一张\emph{虚拟端口}\ 列表，
\eqref{eq:affine-lossless} 中的端口和写成
\begin{equation}
    \sum_{v=1}^{N_v}Y_v(\omega)\,\bm m_v\bm m_v^{\transpose},\qquad
    N_v=\sum_p N_{m,p},
\end{equation}
其中索引 $v$ 与 \code{(projectPortIndex, modeIndex)} 一一对应。
单模 NPM/APM 路径下 $N_{m,p}=1$，$N_v=N_p$ 直接退化。
TFE 多模情况下，每个物理端口可能贡献若干虚拟端口；最多一个被标记为
\code{isExcitationMode}，承载入射幅相，其余作纯吸收端口（在 \eqref{eq:weak} 中
入射项 $\bm E^{\text{inc}}$ 为零）。

\section{\texorpdfstring{对应到 \code{buildAffineSystem()} 的实现}{对应到 buildAffineSystem() 的实现}}
\label{sec:impl}

接口契约见 \code{include/bpfem/fem/FEMAssembler.hpp} 的 \code{AffineSystem}：
\begin{lstlisting}
struct AffineSystem {
    SparseMatrix K;                                               // sum_e (1/mu_r,e) * K_e
    SparseMatrix M;                                               // sum_e (eps_r,e) * M_e
    std::vector<std::vector<std::pair<int,double>>> portCoupling; // m_p (per virtual port)
    std::vector<double> portCutoffSquared;                        // k_{c,p}^2
    std::vector<int>    portFaceIds;                              // physical face id
    std::vector<int>    projectPortIndex;                         // index into project_.ports
    std::vector<bool>   isExcitationMode;                         // true = incident-bearing mode
    std::unordered_set<int> constrainedDofs;                      // PEC zero-Dirichlet mask
    bool lossless = true;                                         // false if any sigma > 0
};
\end{lstlisting}
下面分四步把数学量映射到代码（行号引自 \code{src/fem/FEMAssembler.cpp}）。

\subsection{\texorpdfstring{第一步：单元参考矩阵 $\bm K^T,\bm M^T$ 缓存}{第一步：单元参考矩阵 K\^{}T, M\^{}T 缓存}}
\label{sec:step-cache}

\code{ensureElementCache()} 对每个有效四面体 $T$ 计算一次几何 + 基函数组合
\begin{equation}
    \bigl(\bm K^T_{ab},\bm M^T_{ab}\bigr)
    =\biggl(\sum_q w_q V_T(\nabla\!\times\!\Ne_a)\!\cdot\!(\nabla\!\times\!\Ne_b),\;
    \sum_q w_q V_T \Ne_a\!\cdot\!\Ne_b\biggr)
\end{equation}
其中 $\{(w_q,\bm\lambda_q)\}$ 是 \code{tetraQuadraturePoints()} 给出的 11 点
四面体重心坐标求积公式。对应代码：
\begin{lstlisting}
for (a = 0..dof) for (b = a..dof)
  for each quadrature point qp:
    curlCurl += qp.weight * volume * dot(rowBasis.curl, colBasis.curl);
    mass     += qp.weight * volume * dot(rowBasis.value, colBasis.value);
entry.curlCurlGeo[idx] = curlCurl;        // K^T_{ab}
entry.massGeo    [idx] = mass;            // M^T_{ab}
entry.invMu      = 1 / mu_r,T;
entry.epsR       = eps_r,T;
entry.conductivity = sigma_T;             // used for lossless check
\end{lstlisting}
$\bm K^T,\bm M^T$ 用上三角紧致索引存储 \code{upperTriIndex(a,b,dof)}，
保证装配后 $\bm K,\bm M$ 都是 CSR 上三角对称矩阵。

\subsection{\texorpdfstring{第二步：稀疏图 + 体积装配（$\bm K,\bm M$）}{第二步：稀疏图 + 体积装配（K, M）}}
\label{sec:step-vol}

\code{buildAffineSystem()} 不调用与 \code{assemble()} 共享的稀疏图，而是单独构造
一个\emph{只含体积耦合}\ 的稀疏图：
\begin{lstlisting}
const auto pattern = volumePatternFor(out.constrainedDofs);
out.K = assembleVolumePiece(pattern, out.constrainedDofs, /*useCurlCurl=*/true);
out.M = assembleVolumePiece(pattern, out.constrainedDofs, /*useCurlCurl=*/false);
\end{lstlisting}
理由：ALPS 离线保存的是体积矩阵 $\bm K,\bm M$ 加若干低秩端口向量 $\bm m_v$，
而不是把 $\bm m_v\bm m_v^{\transpose}$ 预先写进稀疏矩阵。因此 $\bm K,\bm M$ 的稀疏图
\emph{严格小于等于}\ 单频点 direct path 的图；端口外积在 ROM 投影或 reduced evaluation 中按低秩形式单独处理。
\code{assembleVolumePiece} 内对每个单元做
\begin{equation}
    \bm K\mathrel{+}=\frac{1}{\mu_{r,T}}\bm K^T,\qquad
    \bm M\mathrel{+}=\varepsilon_{r,T}\bm M^T,
\end{equation}
其中 $\bm K^T,\bm M^T$ 来自第一步缓存，$\sigma_T$ 项故意\emph{不}\ 进入 $\bm M$，
保留实矩阵性质（满足 \eqref{eq:K-M-global}）。

\subsection{第三步：损耗标志}
\label{sec:lossless-flag}

\begin{lstlisting}
out.lossless = true;
for (const auto& entry : elementCache_) {
    if (entry.valid && entry.conductivity > 0.0) {
        out.lossless = false;
        break;
    }
}
\end{lstlisting}
对应数学事实：单一实对称 $\bm M$ 不再能表达 $\varepsilon_c(\omega)$ 在 $\sigma>0$ 时的
$\jj/\omega$ 频率依赖，因此 \eqref{eq:affine-lossless} 不再充分。\code{AlpsSweep}
读到 \code{lossless==false} 时直接抛 \code{runtime\_error} 提示 \code{--sweep direct}。

\subsection{有损介质未仿射化}
\label{sec:lossy-nonaffine}

当计算域中存在 $\sigma_T>0$ 的单元时，\code{buildAffineSystem()} 目前不会继续尝试把有损项写成完整的仿射和，而是直接把 \code{out.lossless} 置为 \code{false}。从连续方程看，问题出在复介电常数
\begin{equation}
    \varepsilon_c(\omega)=\varepsilon_r-\jj\frac{\sigma}{\omega\varepsilon_0}
\end{equation}
所引入的频率依赖并非单一的 $k_0^2(\omega)$ 权重，而是同时包含一个 $\omega^{-1}$ 因子。于是体积项在弱式中可分解为
\begin{equation}
    -k_0^2(\omega)\int_\Omega \varepsilon_r\,\bm v\cdot\bm E\,\dd V
    +\jj\frac{k_0^2(\omega)}{\omega\varepsilon_0}
    \int_\Omega \sigma\,\bm v\cdot\bm E\,\dd V.
\end{equation}
若希望继续保持“离线缓存固定矩阵、在线仅做标量加权”的仿射框架，则除 $\bm K$ 与 $\bm M$ 外，还必须额外引入导电质量矩阵
\begin{equation}
    \bm M_\sigma=\sum_T \sigma_T\bm M^T,
\end{equation}
才能把系统写成
\begin{equation}
    \bm A(\omega)=\bm K-k_0^2(\omega)\bm M+\jj\frac{k_0^2(\omega)}{\omega\varepsilon_0}\bm M_\sigma
    +\sum_vY_v(\omega)\bm m_v\bm m_v^{\transpose}.
    \label{eq:lossy-affine-expanded}
\end{equation}
但这还只是“形式上可写”，并不能直接等价于当前 ALPS 实现可消费的标准仿射结构，因为在线递推和 ROM 评估会同时遇到两个问题：
\begin{enumerate}[leftmargin=2em]
\item $\bm M_\sigma$ 带来的系数是 $\jj\,k_0^2(\omega)/(\omega\varepsilon_0)$，它不是简单的多项式函数，也没有被当前 Krylov 递推显式匹配；
\item 端口项与体积有损项耦合后，若要保持严格的矩匹配或有理近似，还需要对展开点附近的频率导数进行额外处理。
\end{enumerate}
因此，当前实现选择更保守的策略：只要存在任何 $\sigma_T>0$，就不再给出“可直接用于 ROM 的 lossless affine system”，而是把全空间求解交给 \code{assembler.assemble(omega)} 的逐频点 direct path。

从工程实现角度，这个判定等价于：离线阶段只缓存 $\bm K$ 和 $\bm M$ 这两块频率无关矩阵；当有损介质出现时，若不额外实现 $\bm M_\sigma$ 的缓存、在线重组和相应的 ROM 递推更新，继续使用原先的 \code{AlpsSweep} 低维基会导致幅相误差在频带内累积，尤其是在导电损耗较强、共振峰较尖锐时更明显。故而 \code{lossless==false} 并不是“无法求解”，而是“当前这套仿射 ROM 路径未覆盖该物理情形”的显式标记。

\subsection{第四步：虚拟端口展开}
\label{sec:step-port}

最后由注册的 \code{IBoundaryCondition} 列表（典型为 \code{WavePortBC}）通过
\code{affineContributions(ctx)} 给出每个虚拟端口的 $(\bm m_v,k_{c,v}^{2},\dots)$。
\code{WavePortBC::affineContributions} 内部对每个 \code{ProjectPort} 先看
\code{portModeSolver.multiMode(faceId)}：
\begin{enumerate}
\item TFE 多模（\code{multiMode != nullptr}）：按 \code{modeIndex} 顺序把每个解析模 $m$
都打成一个 \code{AffinePortContribution}，
\code{isExcitationMode} 仅在 $m=\,$\code{excitationModeIndex} 且 \code{port.excited}
同时成立时为真；
\item NPM/APM 单模（\code{multiMode == nullptr}）：直接调用
\code{portModeSolver.solve(faceId)}（NPM 走二维 H(curl) 本征，APM 走闭式 TE10 投影），
产出唯一一个虚拟端口。
\end{enumerate}
对每个 contribution 内部还会过滤掉落在 \code{ctx.constrainedDofs} 里的 dof，
保证 $\bm m_v$ 与 $\bm K,\bm M$ 共用同一全局 dof 编号空间。汇总循环是：
\begin{lstlisting}
for (const auto& cond : bcs_) {
    if (!cond) continue;
    for (auto& contrib : cond->affineContributions(ctx)) {
        out.portFaceIds.push_back     (contrib.faceId);
        out.portCutoffSquared.push_back(contrib.cutoffSquared);
        out.portCoupling.push_back    (std::move(contrib.coupling));
        out.projectPortIndex.push_back(contrib.projectPortIndex);
        out.isExcitationMode.push_back(contrib.isExcitationMode);
    }
}
\end{lstlisting}

\bigskip

把这四步串起来，\code{buildAffineSystem()} 实际给出的就是 \eqref{eq:affine-lossless}
所需要的矩阵与端口素材；右端 $\bm b(\omega)$ 不直接缓存，而是由
\code{projectPortIndex}、\code{isExcitationMode}、\code{ProjectPort::magnitudeW/phaseDeg}
以及 \code{powerNormalizationFactor} 在 \code{AlpsSweep::evaluate} 中重建：

\begin{table}[h]
\centering
{\small
\setlength{\tabcolsep}{4pt}
\begin{tabular}{>{\raggedright\arraybackslash}p{0.30\textwidth}
                >{\raggedright\arraybackslash}p{0.25\textwidth}
                >{\raggedright\arraybackslash}p{0.34\textwidth}}
\toprule
数学量 & 代码字段 & 装配位置 \\ \midrule
$\bm K=\sum_T\mu_{r,T}^{-1}\bm K^T$       & \code{out.K}                & §\ref{sec:step-vol} \code{assembleVolumePiece(...,true)}  \\
$\bm M=\sum_T\varepsilon_{r,T}\bm M^T$    & \code{out.M}                & §\ref{sec:step-vol} \code{assembleVolumePiece(...,false)} \\
$\bm m_v$（每虚拟端口）                    & \code{out.portCoupling[v]}  & §\ref{sec:step-port} \code{affineContributions} \\
$k_{c,v}^{2}$                             & \code{out.portCutoffSquared[v]} & 同上 \\
$v\to$ 物理端口面 id                       & \code{out.portFaceIds[v]}   & 同上 \\
$v\to$ ProjectPort 索引                    & \code{out.projectPortIndex[v]} & 同上 \\
是否激励模                                  & \code{out.isExcitationMode[v]} & 同上 \\
入射幅相 $a_v^{\rm inc}$                    & \code{project.ports[...]} & \code{AlpsSweep::evaluate} \\
PEC 约束 dof 集                             & \code{out.constrainedDofs}  & §\ref{sec:pec} \code{cachedConstrainedDofs} \\
$\sigma\equiv0$ 判定                       & \code{out.lossless}         & §\ref{sec:lossless-flag} \\
\bottomrule
\end{tabular}
}
\caption{\code{AffineSystem} 字段与数学模型量的一一对应。}
\end{table}

\section{PEC 零 Dirichlet 投影}
\label{sec:pec}

PEC 面 $\Gamma_e$ 上要求 $\bm n\times\bm E=0$，棱元离散后即"$\Gamma_e$ 上所有切向边 dof 强置零"。
代码用 \code{collectPecConstrainedDofs()} 扫一次表面三角列表，把不属于任何端口的边
全部塞进 \code{cachedConstrainedDofs}。这个集合在 \code{buildAffineSystem()} 里有三处用途：
\begin{enumerate}[leftmargin=2em]
\item \textbf{稀疏图}：\code{volumePatternFor} 中跳过任意端含约束 dof 的物理耦合 $(a,b)$ 对，
但为每个约束 dof 保留一个对角占位，保证矩阵维度仍等于全局边 dof 数。
\item \textbf{累加}：\code{assembleVolumePiece} 中重复同样的过滤，单元贡献不写入约束行/列。
\item \textbf{端口耦合}：\code{WavePortBC::affineContributions} 中对每个 $\bm m_v$ 内部
过滤约束 dof，避免端口耦合向量的支撑超出 $\bm K,\bm M$ 的稀疏图。
\end{enumerate}

注意 \code{buildAffineSystem()} \emph{不}\ 在 $\bm K,\bm M$ 上做 \code{imposeZeroDirichlet}
那样的"主对角线置 1"操作；它选择的是\emph{投影}\ 路径——所有涉及约束 dof 的物理贡献均为零，
只保留对角占位以维持全局维度。数学上这等价于在子空间
$\{x_i=0,\,i\in\text{PEC}\}$ 上写下系统；约束 dof 集仍随 \code{AffineSystem}
传给 \code{AlpsSweep}，方便下游做一致性检查或重建满空间向量。离线阶段用于生成 Krylov 基的
$\bm A_0$ 仍来自 \code{assembler.assemble(...)}，其中 direct path 会对这些对角占位施加单位对角。

\section{装配后的不变量}
\label{sec:invariants}

\begin{enumerate}[leftmargin=2em]
\item \textbf{$\bm K,\bm M$ 实对称}：上三角 CSR 复矩阵但虚部恒为 0，由 \eqref{eq:K-M-global}
本身的实数性保证；引入张量介质后需要扩展为多分量缓存。
\item \textbf{虚拟端口排序}：\code{out.portCoupling} 等数组按
"\code{projectPortIndex} 升序、其内 \code{modeIndex} 升序" 排列，与
\code{WavePortBC::affineContributions} 的循环次序一致。\code{AlpsSweep} 中的
\code{dominantVirtualByProject[]} 依赖这个顺序定位每个物理端口的主模虚拟端口。
\item \textbf{激励唯一性}：每个 \code{projectPortIndex} 至多一个 \code{isExcitationMode==true}。
NPM/APM 单模情形下"唯一" 等价于"全部"。
\item \textbf{$\bm m_v$ 与 $\bm K,\bm M$ 同 dof 空间}：约束 dof 在三处一致剔除，
所以 $\bm m_v$ 不会引用被投影掉的未知量；但 $\bm m_v\bm m_v^{\transpose}$ 一般会在端口 dof 间形成
比体积矩阵更稠密的外积，因此不能要求它与 $\bm K,\bm M$ 同稀疏图。当前 affine path 正是把这些端口项保存为低秩向量而非稀疏矩阵。
\item \textbf{复对称投影}：下游 ALPS 用不带共轭的双线性 $\bm V^{\transpose}\bm K\bm V$、
$\bm V^{\transpose}\bm M\bm V$ 与 $(\bm V^{\transpose}\bm m_v)(\bm V^{\transpose}\bm m_v)^{\transpose}$ 做
Galerkin 投影。$\bm K,\bm M$ 的实对称性 $+$ 端口低秩外积的复对称性
共同保证 reduced matrix 在 \eqref{eq:affine-lossless} 下保留复对称结构。
\end{enumerate}

\section{\texorpdfstring{下游用法：构造 \code{AlpsSweep} 的 ROM 矩阵}{下游用法：构造 AlpsSweep 的 ROM 矩阵}}
\label{sec:alps-downstream}

为完整说明 \code{buildAffineSystem()} 输出的"够用性"，简述下游 \code{AlpsSweep::buildOffline}
对它的消费方式。与其把这一过程描述成普通代码流程，不如把它写成更接近论文中的算法环境，以便突出“离线构基—在线评估”的两阶段结构。

\begin{algorithm}[h]
\caption{由仿射系统构造 \code{AlpsSweep} 的 ROM}
\label{alg:alps-rom}
\begin{algorithmic}[1]
\Require 仿射数据 $\mathcal{A}_{\mathrm{aff}}=(\bm K,\bm M,\{\bm m_v\},\{k_{c,v}^2\},\{a_v^{\rm inc}\})$, 展开频率 $\omega_0$, 阶数 $r$
\Ensure ROM 基 $\bm V$ 及小矩阵 $\bm K_r,\bm M_r,\{\bm m_{v,r}\}$
\State 组装并因式分解 $\bm A_0 \gets \bm A(\omega_0)$
\For{each virtual port $v$}
    \State Solve $\bm A_0\bm w_v^{(1)}=\bm m_v$ \Comment{shift-and-invert 初始向量}
    \For{$k=1$ to $r-1$}
        \State Solve $\bm A_0\bm w_v^{(k+1)}=\bm M\bm w_v^{(k)}$
        \State Orthonormalize $\bm w_v^{(k+1)}$ against existing basis vectors
    \EndFor
\EndFor
\State Assemble orthonormal basis $\bm V=[\bm v_1,\bm v_2,\dots,\bm v_r]$
\State Project $\bm K_r\gets\bm V^{\transpose}\bm K\bm V$, $\bm M_r\gets\bm V^{\transpose}\bm M\bm V$
\For{each virtual port $v$}
    \State Project $\bm m_{v,r}\gets\bm V^{\transpose}\bm m_v$
\EndFor
\State \Return $\bm V,\bm K_r,\bm M_r,\{\bm m_{v,r}\}$
\end{algorithmic}
\end{algorithm}

在中心频率 $\omega_0$ 处，\code{AlpsSweep} 先调用 \code{assembler.assemble(omega0)} 得到与 direct path 完全一致的全空间矩阵 $\bm A_0$ 并做因式分解。
在理想传播模 affine 记号下，它对应
\begin{equation}
    \bm A_0\approx\bm K-k_0^{2}(\omega_0)\bm M
    +\sum_vY_v(\omega_0)\,\bm m_v\bm m_v^{\transpose}.
\end{equation}
这里写成“约等于”是因为 direct path 在低于截止模上有 \code{j*k0} 占位导纳，而当前 reduced evaluation 只重建 $Y_v=\jj\beta_v^+$ 的传播部分。接着对每个虚拟端口求解 $\bm A_0\bm w_v^{(1)}=\bm m_v$ 作为 shift-and-invert Krylov 子空间起始向量；递推
\begin{equation}
    \bm A_0\bm w_v^{(k+1)}=\bm M\bm w_v^{(k)}
\end{equation}
（注意右端只用到\emph{频率无关} 的 $\bm M$）生成 $r$ 阶子空间，再做 modified Gram--Schmidt 正交化得到 $\bm V$。
因此当前实现更准确地说是\emph{单点、M-dominant 的 shift-and-invert Krylov--Galerkin ROM}：它复用了 ALPS/矩匹配文献中“一次或少数几次全空间因式分解 + 小型有理模型”的思想 \cite{Bracken1998,Sun2001,Anderson2004}，但没有显式构造 Lanczos 三对角矩阵，也没有对 $\beta_v(\omega)$ 的导数项做严格 Padé 矩匹配。
投影后的 ROM
\begin{equation}
    \bm K_r=\bm V^{\transpose}\bm K\bm V,\qquad
    \bm M_r=\bm V^{\transpose}\bm M\bm V,\qquad
    \bm m_{v,r}=\bm V^{\transpose}\bm m_v
\end{equation}
是几十到几百维的稠密小矩阵，最终在每个频点 $\omega$ 求解
\begin{equation}
    \Bigl(\bm K_r-k_0^{2}\bm M_r+\sum_vY_v(\omega)\bm m_{v,r}\bm m_{v,r}^{\transpose}\Bigr)
    \bm x_r(\omega)=\bm b_r(\omega)
\end{equation}
其中 $\bm b_r(\omega)=2\sum_vY_v(\omega)a_v^{\rm inc}\bm m_{v,r}$，仅激励虚拟端口贡献非零项。
即可获得端口响应，避免重建全空间场。整条链路对 $\bm K,\bm M,\bm m_v,k_{c,v}^{2}$ 的需求都由 \code{buildAffineSystem()} 一次性满足。

\section{当前限制与扩展点}
\label{sec:limits}

\begin{enumerate}[leftmargin=2em]
\item \textbf{有损介质未仿射化}：$\sigma>0$ 时 \code{lossless==false}，调用方走 direct path。
延伸方向是缓存 $\bm M_\sigma=\sum_T\sigma_T\bm M^T$ 作为第三块全局矩阵，
仿射展开多一项 $\jj\,(k_0^{2}/(\omega\varepsilon_0))\bm M_\sigma$。若要做严格矩匹配，
Krylov 递推还需要把该项对展开变量的导数纳入；当前实现并未这样做。
\item \textbf{各向异性介质}：\eqref{eq:Ke}--\eqref{eq:Me} 假设 $\mu_r,\varepsilon_r$ 标量，
对张量介质需要把 $\bm K^T,\bm M^T$ 拆成多分量并各自加权。
\item \textbf{端口模式频率无关性}：\eqref{eq:m_p} 假设 $\bm e_{p,m}^{\perp}$ 与 $\omega$ 无关，
这是 ALPS 仿射展开的代数前提。NPM/TFE 路径下这要求二维 H(curl) 本征问题
\eqref{eq:port-eig} 的特征对（在频带内不发生模式交叉的前提下）频率无关，
对空气或常规均匀填充波导近似成立；色散填充、强损耗端口或开放波导需要单独处理。
APM 路径假设矩形 TE10 闭式同样频率无关。
\item \textbf{低于截止模}：direct path 当前用 \code{j*k0} 作为 $\beta^+=0$ 时的占位导纳，
而 ALPS reduced evaluation 只在 $\beta^+>0$ 时加入端口低秩项和入射右端。若 TFE 保留了频带内
低于截止的高阶模，direct/ALPS 两条路径会在该部分端口算子上不完全一致；严格模型应使用完整的
复传播常数或分传播/衰减两类模分别仿射化。
\item \textbf{ABC/SIBC 边界条件}：当前 \code{IBoundaryCondition::affineContributions} 默认空。
吸收边界、阻抗边界一般给出多项形式（含 $\sqrt{\jj\omega}$ 等），
若要纳入仿射展开需要额外的项。
\end{enumerate}

\section{小结}

\code{FEMAssembler::buildAffineSystem()} 把频域 curl--curl FEM 系统在 lossless 假设下
彻底拆成 \eqref{eq:affine-lossless} 的"频率无关矩阵 $+$ 频率函数加权"形式，
为 ALPS/Krylov 类快速扫频提供可复用的代数基础。其实现遵循三条主线：
单元参考矩阵缓存（§\ref{sec:step-cache}）、体积全局矩阵装配（§\ref{sec:step-vol}）、
经 \code{IBoundaryCondition} 分发的虚拟端口展开（§\ref{sec:step-port}）；
PEC 投影（§\ref{sec:pec}）和损耗标志（§\ref{sec:lossless-flag}）作为两条横切机制
保证后续 ROM 与单频点 direct path 的等价性。

\begin{thebibliography}{9}
\bibitem{Nedelec1980}
J. C. Nédélec,
``Mixed finite elements in $\mathbb{R}^3$,''
\emph{Numerische Mathematik}, vol. 35, pp. 315--341, 1980.

\bibitem{Graglia1997}
R. D. Graglia, D. R. Wilton, and A. F. Peterson,
``Higher order interpolatory vector bases for computational electromagnetics,''
\emph{IEEE Transactions on Antennas and Propagation}, vol. 45, no. 3, pp. 329--342, 1997.

\bibitem{Webb1999}
J. P. Webb,
``Hierarchal vector basis functions of arbitrary order for triangular and tetrahedral finite elements,''
\emph{IEEE Transactions on Antennas and Propagation}, vol. 47, no. 8, pp. 1244--1253, 1999.

\bibitem{Cendes1988}
Z. J. Cendes and J.-F. Lee,
``The transfinite element method for modeling MMIC devices,''
\emph{IEEE Transactions on Microwave Theory and Techniques}, vol. 36, no. 12, pp. 1639--1649, 1988.

\bibitem{Lee1990}
J.-F. Lee,
``Analysis of passive microwave devices by using three-dimensional tangential vector finite elements,''
\emph{International Journal of Numerical Modelling}, vol. 3, no. 4, pp. 235--246, 1990.

\bibitem{Bracken1998}
J. E. Bracken, D.-K. Sun, and Z. J. Cendes,
``S-domain methods for simultaneous time and frequency characterization of electromagnetic devices,''
\emph{IEEE Transactions on Microwave Theory and Techniques}, vol. 46, no. 9, pp. 1277--1290, 1998.

\bibitem{Sun2001}
D.-K. Sun, Z. Cendes, and J.-F. Lee,
``ALPS: A new fast frequency-sweep procedure for microwave devices,''
\emph{IEEE Transactions on Microwave Theory and Techniques}, vol. 49, no. 2, pp. 398--402, 2001.

\bibitem{Anderson2004}
B. Anderson, J. E. Bracken, J. B. Manges, G. Peng, and Z. Cendes,
``Full-wave analysis in SPICE via model-order reduction,''
\emph{IEEE Transactions on Microwave Theory and Techniques}, vol. 52, no. 9, pp. 2314--2320, 2004.
\end{thebibliography}

\section{\texorpdfstring{下游用法：构造 \code{AlpsSweep} 的 ROM 矩阵}{下游用法：构造 AlpsSweep 的 ROM 矩阵}}
\label{sec:downstream-rom}

这一节专门回答一个问题：\code{buildAffineSystem()} 产出的 $\bm K,\bm M,\bm m_v,k_{c,v}^2$
为什么足以在 \code{AlpsSweep} 中构造 ROM 矩阵。结合 ALPS 原始论文中“只做一次全空间分解、在线反复评估低维模型”的思路
\cite{Sun2001,Bracken1998,Anderson2004}，这里可以把整个过程分成三层：
(1) 频率无关的离线数据如何定义；
(2) 这些数据如何通过 Krylov / Galerkin 投影生成小矩阵；
(3) 小矩阵如何重建每个频点的端口响应。

\subsection{从全空间仿射系统到频点矩阵}

设 \code{buildAffineSystem()} 给出离线对象
\begin{equation}
    \mathcal{A}_{\mathrm{aff}}=\bigl(\bm K,\bm M,\{\bm m_v\}_{v=1}^{N_v},\{k_{c,v}^2\}_{v=1}^{N_v},\{a_v^{\rm inc}\}_{v=1}^{N_v}\bigr).
\end{equation}
在 lossless 且端口模频率独立的前提下，任意频点的全空间矩阵都可写成
\begin{equation}
    \bm A(\omega)=\bm K-k_0^2(\omega)\bm M+\sum_{v=1}^{N_v}Y_v(\omega)\bm m_v\bm m_v^{\transpose},
    \qquad Y_v(\omega)=\jj\beta_v^+(\omega).
    \label{eq:alps-full-affine}
\end{equation}
如果存在入射激励，则右端满足
\begin{equation}
    \bm b(\omega)=2\sum_{v\in\mathcal{E}}Y_v(\omega)a_v^{\rm inc}\bm m_v,
\end{equation}
其中 $\mathcal{E}$ 是被标记为激励模的虚拟端口集合。式 \eqref{eq:alps-full-affine} 的关键意义是：
频率依赖只出现在标量函数 $k_0^2(\omega)$ 与 $Y_v(\omega)$ 中，矩阵结构本身是固定的。
这正是 \code{AlpsSweep} 能把“全矩阵扫频”转化成“少量小矩阵扫频”的根本原因。

\subsection{离线阶段：以中心频率建立展开基}

\code{AlpsSweep::buildOffline()} 首先选一个展开中心频率 $\omega_0$，并调用
\code{assembler.assemble(omega0)} 得到与 direct path 一致的全空间矩阵
\begin{equation}
    \bm A_0 := \bm A(\omega_0).
\end{equation}
在理想传播模记号下，$\bm A_0$ 既包含体积项 $\bm K-k_0^2(\omega_0)\bm M$，也包含中心频率处的端口秩 1 项。
随后对每个虚拟端口 $v$ 解线性方程
\begin{equation}
    \bm A_0\bm w_v^{(1)}=\bm m_v,
    \label{eq:start-vector}
\end{equation}
这一步是 ALPS / AWE / Krylov 族方法中典型的 shift-and-invert 起步：
把“端口耦合向量”经过 $\bm A_0^{-1}$ 滤波后，得到最能反映该端口局部频响的初始方向。
如果某个端口有多个保留模，则每个虚拟端口各自产生一个起始向量。

接下来，\code{AlpsSweep} 不再直接使用原始的频域矩阵，而是用固定的递推关系扩展子空间。当前实现的核心递推是
\begin{equation}
    \bm A_0\bm w_v^{(k+1)}=\bm M\bm w_v^{(k)},\qquad k=1,2,\dots,r-1.
    \label{eq:krylov-recursion}
\end{equation}
这里右端只用到频率无关的 $\bm M$，因此每次递推都等价于“先乘一次质量矩阵，再用同一个分解好的 $\bm A_0$ 回代”。
这就把大量频点的复杂性压缩进了一个公共子空间中。

为了避免不同端口、不同阶次的向量线性相关，递推得到的所有向量会做 modified Gram--Schmidt 正交化，形成正交基矩阵
\begin{equation}
    \bm V=[\bm v_1,\bm v_2,\dots,\bm v_r],\qquad \bm V^{\transpose}\bm V=\bm I.
\end{equation}
这一步在实现上可以概括为：对每个虚拟端口分别启动一次 shift-and-invert 递推，逐阶生成向量，再把全体向量汇总后做正交化与去重。其算法结构可写成如下伪代码。

\begin{lstlisting}[language=C++,caption={\code{AlpsSweep} 离线阶段的虚拟端口子空间构造伪代码},label={lst:alps-offline-pseudocode}]
Input:
    A0 = assemble(omega0)        // center-frequency full matrix
    M                              // frequency-independent mass matrix
    virtualPorts = {m_v}
    r                              // target subspace order per start vector
Output:
    V                              // orthonormal basis for ROM

W = []
for each virtual port v in virtualPorts:
    w1 = solve(A0, m_v)
    append w1 to W

    w_prev = w1
    for k = 1 .. r-1:
        rhs = M * w_prev
        w_next = solve(A0, rhs)
        append w_next to W
        w_prev = w_next

V = []
for each w in W:
    for each q in V:
        w = w - q * (q^T * w)      // modified Gram-Schmidt, 1st pass
    for each q in V:
        correction = q * (q^T * w) // optional reorthogonalization pass
        w = w - correction
    if norm(w) > tol:
        w = w / norm(w)
        append w to V
return V
\end{lstlisting}

于是，\code{AlpsSweep} 的离线阶段本质上是在构造一个“对全空间响应最有代表性”的低维试探空间。

\subsection{ROM 矩阵是如何投影出来的}

一旦有了基 \(\bm V\)，ROM 矩阵就通过 Galerkin 投影得到：
\begin{equation}
    \bm K_r=\bm V^{\transpose}\bm K\bm V,\qquad
    \bm M_r=\bm V^{\transpose}\bm M\bm V,\qquad
    \bm m_{v,r}=\bm V^{\transpose}\bm m_v.
    \label{eq:rom-projection}
\end{equation}
为什么这样做是合理的？因为原系统的每一项都满足“左乘测试空间、右乘试探空间”的标准有限元变分结构：
体积项是二次型 $\bm x^{\transpose}\bm K\bm x$ 与 $\bm x^{\transpose}\bm M\bm x$；端口项是低秩外积
$\bm m_v\bm m_v^{\transpose}$。在 Galerkin 框架下，只要把未知场近似为
\begin{equation}
    \bm x(\omega)\approx\bm V\bm x_r(\omega),
\end{equation}
并用同一个 $\bm V$ 左乘残差，就会自然得到小系统
\begin{equation}
    \bm V^{\transpose}\bm A(\omega)\bm V\,\bm x_r(\omega)=\bm V^{\transpose}\bm b(\omega).
\end{equation}
将 \eqref{eq:alps-full-affine} 代入后，立刻得到
\begin{equation}
    \Bigl(\bm K_r-k_0^2(\omega)\bm M_r+\sum_{v=1}^{N_v}Y_v(\omega)\bm m_{v,r}\bm m_{v,r}^{\transpose}\Bigr)
    \bm x_r(\omega)=\bm b_r(\omega),
    \label{eq:rom-system}
\end{equation}
其中
\begin{equation}
    \bm b_r(\omega)=2\sum_{v\in\mathcal{E}}Y_v(\omega)a_v^{\rm inc}\bm m_{v,r}.
\end{equation}
这说明 ROM 并不是“重新拟合”全系统，而是把同一仿射结构投影到了低维空间。

\subsection{为什么端口项在 ROM 中仍保持秩 1}

端口项在全空间是 $\bm m_v\bm m_v^{\transpose}$，它的秩至多为 1。投影后：
\begin{equation}
    \bm V^{\transpose}(\bm m_v\bm m_v^{\transpose})\bm V
    =(\bm V^{\transpose}\bm m_v)(\bm V^{\transpose}\bm m_v)^{\transpose}
    =\bm m_{v,r}\bm m_{v,r}^{\transpose}.
\end{equation}
所以端口项的低秩结构被完全保留，只是向量长度从全空间维数 $N$ 变成了 ROM 维数 $r$。
这也是 \code{AlpsSweep} 的关键优势之一：频率变化时，只需更新标量系数 $Y_v(\omega)$，无需重装配稀疏矩阵。

\subsection{和 ALPS 文献中的矩阵观点对应起来}

从文献角度看，\cite{Bracken1998,Sun2001,Anderson2004} 中的 ALPS / AWE / MOR 方法，
本质上都是把一个关于频率的矩阵函数写成“若干固定矩阵 + 标量函数”的形式，再围绕某个展开点构造低维近似。
这里的 \code{AlpsSweep} 做法与其一致，只是实现上更贴近有限元工程：
\begin{enumerate}[leftmargin=2em]
\item 先把有限元问题拆成频率无关的 \(\bm K,\bm M,\bm m_v\)；
\item 再用 \(\bm A_0^{-1}\) 驱动 Krylov 递推生成基；
\item 最后在小空间里重建每个频点的矩阵 \eqref{eq:rom-system}。
\end{enumerate}
因此，所谓“构造 ROM 矩阵”，并不是凭空造一个新模型，而是对全空间仿射系统做严格的投影降维。

\subsection{为何 \code{buildAffineSystem()} 足够}

注意到 \eqref{eq:rom-system} 里所有需要的数据都已经由 \code{buildAffineSystem()} 提供：
\begin{itemize}[leftmargin=2em]
\item \(\bm K\) 与 \(\bm M\)：决定体积刚度与质量；
\item \(\bm m_v\)：决定每个虚拟端口的秩 1 耦合；
\item \(k_{c,v}^2\)：决定 \(Y_v(\omega)\) 的截止行为；
\item \(a_v^{\rm inc}\) 与 \code{isExcitationMode}：决定右端是否注入入射波。
\end{itemize}
换句话说，\code{AlpsSweep} 构造 ROM 所需的一切代数素材都已经在离线阶段被标准化了；剩下的工作只是“选基 + 投影 + 小矩阵扫频”。

\medskip

因此，\code{buildAffineSystem()} 输出的不是某个单频点矩阵，而是一个可被 \code{AlpsSweep}
直接消费的仿射算子集合。ROM 矩阵的构造公式最终就是 \eqref{eq:rom-projection} 和 \eqref{eq:rom-system}。
在数值上，这让大规模三维端口问题从“每个频点都解一次大稀疏系统”变成“离线一次分解 + 在线反复解小稠密系统”，
这正是 ALPS 的加速来源。

\end{document}
